\documentclass[11pt,a4paper]{article}

% Packages essentiels
\usepackage[utf8]{inputenc}
\usepackage[french]{babel}
\usepackage[T1]{fontenc}
\usepackage{geometry}
\geometry{margin=2cm}
\usepackage{xcolor}
\usepackage{listings}
\usepackage{tcolorbox}
\usepackage{fancyhdr}
\usepackage{titlesec}
\usepackage{graphicx}
\usepackage{enumitem}
\usepackage{hyperref}
\usepackage{amsmath}
\usepackage{amssymb}

% Configuration des couleurs
\definecolor{codebackground}{RGB}{245,245,245}
\definecolor{commentcolor}{RGB}{0,128,0}
\definecolor{keywordcolor}{RGB}{0,0,255}
\definecolor{stringcolor}{RGB}{163,21,21}
\definecolor{functioncolor}{RGB}{128,0,128}
\definecolor{numbercolor}{RGB}{255,127,0}
\definecolor{titlecolor}{RGB}{0,51,102}
\definecolor{sectioncolor}{RGB}{0,102,204}
\definecolor{importantcolor}{RGB}{204,0,0}

% Configuration du style de code
\lstset{
    language=C,
    backgroundcolor=\color{codebackground},
    basicstyle=\ttfamily\small,
    keywordstyle=\color{keywordcolor}\bfseries,
    commentstyle=\color{commentcolor}\itshape,
    stringstyle=\color{stringcolor},
    numberstyle=\tiny\color{numbercolor},
    numbers=left,
    numbersep=8pt,
    frame=single,
    frameround=tttt,
    rulecolor=\color{black!30},
    breaklines=true,
    breakatwhitespace=true,
    tabsize=4,
    showstringspaces=false,
    captionpos=b,
    xleftmargin=15pt,
    xrightmargin=5pt,
    escapeinside={(*@}{@*)},
    morekeywords={pthread_t, pthread_mutex_t, pthread_cond_t, pthread_attr_t}
}

% Style des titres
\titleformat{\section}
{\color{sectioncolor}\Large\bfseries}
{\thesection}{1em}{}[\titlerule]

\titleformat{\subsection}
{\color{sectioncolor}\large\bfseries}
{\thesubsection}{1em}{}

% En-tête et pied de page
\pagestyle{fancy}
\fancyhf{}
\fancyhead[L]{\textbf{Cours Pthreads - Programmation Concurrente}}
\fancyhead[R]{\thepage}
\fancyfoot[C]{\textit{Document de révision pour examen}}

% Boîtes personnalisées
\newtcolorbox{importantbox}[1]{
    colback=red!5!white,
    colframe=importantcolor,
    fonttitle=\bfseries,
    title=#1
}

\newtcolorbox{definitionbox}[1]{
    colback=blue!5!white,
    colframe=sectioncolor,
    fonttitle=\bfseries,
    title=#1
}

\newtcolorbox{examplebox}[1]{
    colback=green!5!white,
    colframe=green!50!black,
    fonttitle=\bfseries,
    title=#1
}

% Titre du document
\title{
    \Huge\textbf{\color{titlecolor}Guide Complet des Pthreads}\\
    \Large\textit{Programmation Concurrente en C}\\
    \vspace{0.5cm}
    \large Document de Révision pour Examen
}
\author{}
\date{\today}

\begin{document}

\maketitle
\thispagestyle{empty}

\newpage
\tableofcontents
\newpage

% ============================================
\section{Introduction aux Pthreads}
% ============================================

\begin{definitionbox}{Qu'est-ce qu'un Thread ?}
Un \textbf{thread} (fil d'exécution) est une unité d'exécution légère au sein d'un processus. Plusieurs threads peuvent s'exécuter simultanément dans le même espace mémoire, partageant les ressources du processus.
\end{definitionbox}

\subsection{Différence entre Processus et Thread}

\begin{center}
\begin{tabular}{|p{5cm}|p{5cm}|}
\hline
\textbf{Processus} & \textbf{Thread} \\
\hline
Espace mémoire séparé & Partage l'espace mémoire \\
\hline
Création coûteuse & Création légère \\
\hline
Communication complexe (IPC) & Communication simple (variables partagées) \\
\hline
Isolation forte & Isolation faible \\
\hline
\end{tabular}
\end{center}

\subsection{Avantages des Pthreads}

\begin{itemize}[leftmargin=*]
    \item \textbf{Performance} : Exploitation du parallélisme sur processeurs multi-cœurs
    \item \textbf{Réactivité} : Interface utilisateur reste responsive pendant calculs lourds
    \item \textbf{Partage de ressources} : Communication efficace entre threads
    \item \textbf{Économie} : Moins coûteux que la création de processus
\end{itemize}

\begin{importantbox}{Compilation avec Pthreads}
Pour compiler un programme utilisant pthreads :
\begin{verbatim}
gcc -pthread programme.c -o programme
\end{verbatim}
L'option \texttt{-pthread} est \textbf{obligatoire} !
\end{importantbox}

% ============================================
\section{Création et Gestion des Threads}
% ============================================

\subsection{Fonction pthread\_create()}

\begin{definitionbox}{Prototype}
\begin{verbatim}
int pthread_create(pthread_t *thread, 
                   const pthread_attr_t *attr,
                   void *(*start_routine)(void *), 
                   void *arg);
\end{verbatim}
\end{definitionbox}

\textbf{Paramètres :}
\begin{itemize}
    \item \texttt{thread} : Pointeur vers l'identifiant du thread créé
    \item \texttt{attr} : Attributs du thread (NULL pour valeurs par défaut)
    \item \texttt{start\_routine} : Fonction que le thread va exécuter
    \item \texttt{arg} : Argument passé à la fonction
\end{itemize}

\textbf{Retour :} 0 en cas de succès, code d'erreur sinon

\subsection{Exemple 1 : Création Simple d'un Thread}

\begin{lstlisting}[caption={Création basique d'un thread}]
#include <stdio.h>          (*@\color{commentcolor}// Entrées/sorties standard@*)
#include <stdlib.h>         (*@\color{commentcolor}// Fonctions utilitaires@*)
#include <pthread.h>        (*@\color{commentcolor}// Bibliothèque pthreads@*)
#include <unistd.h>         (*@\color{commentcolor}// Pour sleep()@*)

(*@\color{commentcolor}// Fonction exécutée par le thread@*)
void* fonction_thread(void* arg) {
    (*@\color{commentcolor}// Conversion de l'argument void* en int*@*)
    int id = *(int*)arg;
    
    (*@\color{commentcolor}// Affichage de l'ID du thread@*)
    printf("Thread %d : Démarrage\n", id);
    
    (*@\color{commentcolor}// Simulation d'un travail (2 secondes)@*)
    sleep(2);
    
    (*@\color{commentcolor}// Message de fin@*)
    printf("Thread %d : Terminaison\n", id);
    
    (*@\color{commentcolor}// Retour NULL (pas de valeur de retour)@*)
    return NULL;
}

int main() {
    (*@\color{commentcolor}// Déclaration de l'identifiant du thread@*)
    pthread_t thread;
    
    (*@\color{commentcolor}// ID à passer au thread@*)
    int id = 1;
    
    (*@\color{commentcolor}// Création du thread@*)
    (*@\color{commentcolor}// - \&thread : où stocker l'ID du thread@*)
    (*@\color{commentcolor}// - NULL : attributs par défaut@*)
    (*@\color{commentcolor}// - fonction\_thread : fonction à exécuter@*)
    (*@\color{commentcolor}// - \&id : argument passé à la fonction@*)
    if (pthread_create(&thread, NULL, fonction_thread, &id) != 0) {
        perror("Erreur pthread_create");
        return 1;
    }
    
    (*@\color{commentcolor}// Le thread principal continue son exécution@*)
    printf("Main : Thread créé avec succès\n");
    
    (*@\color{commentcolor}// Attendre la fin du thread@*)
    pthread_join(thread, NULL);
    
    (*@\color{commentcolor}// Fin du programme@*)
    printf("Main : Programme terminé\n");
    return 0;
}
\end{lstlisting}

\subsection{Fonction pthread\_join()}

\begin{definitionbox}{Prototype}
\begin{verbatim}
int pthread_join(pthread_t thread, void **retval);
\end{verbatim}
\end{definitionbox}

\textbf{Rôle :} Attend la terminaison d'un thread spécifique (équivalent de \texttt{wait()} pour les processus)

\textbf{Paramètres :}
\begin{itemize}
    \item \texttt{thread} : Identifiant du thread à attendre
    \item \texttt{retval} : Pointeur pour récupérer la valeur de retour (NULL si non nécessaire)
\end{itemize}

\subsection{Exemple 2 : Threads Multiples avec Valeurs de Retour}

\begin{lstlisting}[caption={Création de plusieurs threads avec retour de valeurs}]
#include <stdio.h>
#include <stdlib.h>
#include <pthread.h>

(*@\color{commentcolor}// Fonction qui calcule le carré d'un nombre@*)
void* calculer_carre(void* arg) {
    (*@\color{commentcolor}// Récupération du nombre@*)
    int num = *(int*)arg;
    
    (*@\color{commentcolor}// Allocation dynamique pour le résultat@*)
    (*@\color{commentcolor}// (nécessaire car la variable locale disparaît)@*)
    int* resultat = malloc(sizeof(int));
    
    (*@\color{commentcolor}// Calcul du carré@*)
    *resultat = num * num;
    
    printf("Thread : %d^2 = %d\n", num, *resultat);
    
    (*@\color{commentcolor}// Retour du résultat (cast en void*)@*)
    return (void*)resultat;
}

int main() {
    (*@\color{commentcolor}// Tableau de 5 threads@*)
    pthread_t threads[5];
    (*@\color{commentcolor}// Tableau des nombres à traiter@*)
    int nombres[5] = {1, 2, 3, 4, 5};
    
    (*@\color{commentcolor}// Création des 5 threads@*)
    for (int i = 0; i < 5; i++) {
        if (pthread_create(&threads[i], NULL, calculer_carre, &nombres[i]) != 0) {
            fprintf(stderr, "Erreur création thread %d\n", i);
            return 1;
        }
    }
    
    (*@\color{commentcolor}// Attente et récupération des résultats@*)
    for (int i = 0; i < 5; i++) {
        (*@\color{commentcolor}// Pointeur pour récupérer la valeur de retour@*)
        void* retval;
        
        (*@\color{commentcolor}// Attente du thread i@*)
        pthread_join(threads[i], &retval);
        
        (*@\color{commentcolor}// Conversion et affichage du résultat@*)
        int* resultat = (int*)retval;
        printf("Main : Résultat thread %d = %d\n", i, *resultat);
        
        (*@\color{commentcolor}// Libération de la mémoire allouée par le thread@*)
        free(resultat);
    }
    
    return 0;
}
\end{lstlisting}

\subsection{Fonction pthread\_exit()}

\begin{definitionbox}{Prototype}
\begin{verbatim}
void pthread_exit(void *retval);
\end{verbatim}
\end{definitionbox}

\textbf{Rôle :} Termine le thread appelant et retourne une valeur

\begin{importantbox}{Différence return vs pthread\_exit()}
\begin{itemize}
    \item \texttt{return} : Termine uniquement la fonction du thread
    \item \texttt{pthread\_exit()} : Termine le thread immédiatement, même depuis une fonction appelée
\end{itemize}
\end{importantbox}

\subsection{Fonction pthread\_detach()}

\begin{definitionbox}{Prototype}
\begin{verbatim}
int pthread_detach(pthread_t thread);
\end{verbatim}
\end{definitionbox}

\textbf{Rôle :} Détache un thread, permettant la libération automatique de ses ressources à sa terminaison (pas besoin de \texttt{pthread\_join()})

\begin{lstlisting}[caption={Utilisation de pthread\_detach()}]
void* fonction_detachee(void* arg) {
    (*@\color{commentcolor}// Se détacher soi-même@*)
    pthread_detach(pthread_self());
    
    (*@\color{commentcolor}// Travail du thread@*)
    printf("Thread détaché en cours d'exécution\n");
    sleep(1);
    
    return NULL;
}

int main() {
    pthread_t thread;
    pthread_create(&thread, NULL, fonction_detachee, NULL);
    
    (*@\color{commentcolor}// Pas besoin de pthread\_join()@*)
    (*@\color{commentcolor}// Le thread se nettoie automatiquement@*)
    
    sleep(2); (*@\color{commentcolor}// Attendre que le thread finisse@*)
    return 0;
}
\end{lstlisting}

% ============================================
\section{Synchronisation : Mutex}
% ============================================

\begin{definitionbox}{Qu'est-ce qu'un Mutex ?}
Un \textbf{mutex} (Mutual Exclusion) est un verrou qui garantit qu'un seul thread à la fois peut accéder à une section critique (ressource partagée).
\end{definitionbox}

\subsection{Problème de la Section Critique}

\begin{importantbox}{Race Condition}
Quand plusieurs threads accèdent simultanément à une variable partagée sans synchronisation, le résultat est \textbf{imprévisible} !
\end{importantbox}

\subsection{Fonctions Mutex Principales}

\begin{definitionbox}{Déclaration et Initialisation}
\begin{verbatim}
pthread_mutex_t mutex = PTHREAD_MUTEX_INITIALIZER;
// ou
pthread_mutex_t mutex;
pthread_mutex_init(&mutex, NULL);
\end{verbatim}
\end{definitionbox}

\textbf{Fonctions de base :}
\begin{itemize}
    \item \texttt{pthread\_mutex\_lock()} : Verrouille le mutex (attend s'il est déjà verrouillé)
    \item \texttt{pthread\_mutex\_unlock()} : Déverrouille le mutex
    \item \texttt{pthread\_mutex\_trylock()} : Tente de verrouiller sans bloquer
    \item \texttt{pthread\_mutex\_destroy()} : Détruit le mutex
\end{itemize}

\subsection{Exemple 3 : Compteur Partagé SANS Mutex (Problème)}

\begin{lstlisting}[caption={Problème de race condition}]
#include <stdio.h>
#include <pthread.h>

(*@\color{importantcolor}// Variable partagée (DANGER sans protection !)@*)
int compteur = 0;

void* incrementer(void* arg) {
    for (int i = 0; i < 1000000; i++) {
        (*@\color{importantcolor}// SECTION CRITIQUE NON PROTÉGÉE@*)
        (*@\color{importantcolor}// Opération non atomique : lecture, incrémentation, écriture@*)
        compteur++;  (*@\color{importantcolor}// RACE CONDITION !@*)
    }
    return NULL;
}

int main() {
    pthread_t t1, t2;
    
    pthread_create(&t1, NULL, incrementer, NULL);
    pthread_create(&t2, NULL, incrementer, NULL);
    
    pthread_join(t1, NULL);
    pthread_join(t2, NULL);
    
    (*@\color{importantcolor}// Résultat attendu : 2000000@*)
    (*@\color{importantcolor}// Résultat réel : IMPRÉVISIBLE (souvent < 2000000)@*)
    printf("Compteur final : %d\n", compteur);
    
    return 0;
}
\end{lstlisting}

\subsection{Exemple 4 : Compteur Partagé AVEC Mutex (Solution)}

\begin{lstlisting}[caption={Protection avec mutex}]
#include <stdio.h>
#include <pthread.h>

(*@\color{commentcolor}// Variable partagée@*)
int compteur = 0;

(*@\color{commentcolor}// Mutex pour protéger la section critique@*)
pthread_mutex_t mutex = PTHREAD_MUTEX_INITIALIZER;

void* incrementer(void* arg) {
    for (int i = 0; i < 1000000; i++) {
        (*@\color{commentcolor}// VERROUILLAGE : début de la section critique@*)
        pthread_mutex_lock(&mutex);
        
        (*@\color{commentcolor}// Section critique protégée@*)
        (*@\color{commentcolor}// Un seul thread à la fois peut exécuter ce code@*)
        compteur++;
        
        (*@\color{commentcolor}// DÉVERROUILLAGE : fin de la section critique@*)
        pthread_mutex_unlock(&mutex);
    }
    return NULL;
}

int main() {
    pthread_t t1, t2;
    
    pthread_create(&t1, NULL, incrementer, NULL);
    pthread_create(&t2, NULL, incrementer, NULL);
    
    pthread_join(t1, NULL);
    pthread_join(t2, NULL);
    
    (*@\color{commentcolor}// Résultat garanti : 2000000@*)
    printf("Compteur final : %d\n", compteur);
    
    (*@\color{commentcolor}// Destruction du mutex@*)
    pthread_mutex_destroy(&mutex);
    
    return 0;
}
\end{lstlisting}

\subsection{Exemple 5 : Compte Bancaire Partagé}

\begin{lstlisting}[caption={Simulation de transactions bancaires}]
#include <stdio.h>
#include <stdlib.h>
#include <pthread.h>
#include <unistd.h>

(*@\color{commentcolor}// Structure représentant un compte bancaire@*)
typedef struct {
    int solde;                  (*@\color{commentcolor}// Solde du compte@*)
    pthread_mutex_t mutex;      (*@\color{commentcolor}// Mutex pour protéger le solde@*)
} CompteBancaire;

(*@\color{commentcolor}// Structure pour passer les arguments aux threads@*)
typedef struct {
    CompteBancaire* compte;
    int montant;
    int id_thread;
} ArgsTransaction;

(*@\color{commentcolor}// Fonction de dépôt@*)
void* deposer(void* arg) {
    ArgsTransaction* args = (ArgsTransaction*)arg;
    
    (*@\color{commentcolor}// Verrouillage du compte@*)
    pthread_mutex_lock(&args->compte->mutex);
    
    (*@\color{commentcolor}// Lecture du solde actuel@*)
    int solde_avant = args->compte->solde;
    
    (*@\color{commentcolor}// Simulation d'un délai de traitement@*)
    usleep(100);
    
    (*@\color{commentcolor}// Ajout du montant@*)
    args->compte->solde += args->montant;
    
    printf("Thread %d : Dépôt de %d€ (Solde : %d€ -> %d€)\n",
           args->id_thread, args->montant, 
           solde_avant, args->compte->solde);
    
    (*@\color{commentcolor}// Déverrouillage du compte@*)
    pthread_mutex_unlock(&args->compte->mutex);
    
    free(args);
    return NULL;
}

(*@\color{commentcolor}// Fonction de retrait@*)
void* retirer(void* arg) {
    ArgsTransaction* args = (ArgsTransaction*)arg;
    
    pthread_mutex_lock(&args->compte->mutex);
    
    (*@\color{commentcolor}// Vérification du solde suffisant@*)
    if (args->compte->solde >= args->montant) {
        int solde_avant = args->compte->solde;
        usleep(100);
        args->compte->solde -= args->montant;
        
        printf("Thread %d : Retrait de %d€ (Solde : %d€ -> %d€)\n",
               args->id_thread, args->montant,
               solde_avant, args->compte->solde);
    } else {
        printf("Thread %d : Retrait refusé (solde insuffisant)\n",
               args->id_thread);
    }
    
    pthread_mutex_unlock(&args->compte->mutex);
    
    free(args);
    return NULL;
}

int main() {
    (*@\color{commentcolor}// Initialisation du compte avec 1000€@*)
    CompteBancaire compte;
    compte.solde = 1000;
    pthread_mutex_init(&compte.mutex, NULL);
    
    (*@\color{commentcolor}// Création de 10 threads (5 dépôts, 5 retraits)@*)
    pthread_t threads[10];
    
    for (int i = 0; i < 5; i++) {
        (*@\color{commentcolor}// Thread de dépôt@*)
        ArgsTransaction* args_depot = malloc(sizeof(ArgsTransaction));
        args_depot->compte = &compte;
        args_depot->montant = 100;
        args_depot->id_thread = i;
        pthread_create(&threads[i], NULL, deposer, args_depot);
        
        (*@\color{commentcolor}// Thread de retrait@*)
        ArgsTransaction* args_retrait = malloc(sizeof(ArgsTransaction));
        args_retrait->compte = &compte;
        args_retrait->montant = 150;
        args_retrait->id_thread = i + 5;
        pthread_create(&threads[i+5], NULL, retirer, args_retrait);
    }
    
    (*@\color{commentcolor}// Attente de tous les threads@*)
    for (int i = 0; i < 10; i++) {
        pthread_join(threads[i], NULL);
    }
    
    printf("\nSolde final : %d€\n", compte.solde);
    
    (*@\color{commentcolor}// Nettoyage@*)
    pthread_mutex_destroy(&compte.mutex);
    
    return 0;
}
\end{lstlisting}

% ============================================
\section{Synchronisation : Variables de Condition}
% ============================================

\begin{definitionbox}{Qu'est-ce qu'une Variable de Condition ?}
Une \textbf{variable de condition} permet à un thread d'attendre qu'une certaine condition soit vraie avant de continuer. Elle est toujours utilisée avec un mutex.
\end{definitionbox}

\subsection{Fonctions Principales}

\begin{itemize}
    \item \texttt{pthread\_cond\_init()} : Initialise une variable de condition
    \item \texttt{pthread\_cond\_wait()} : Attend un signal (libère le mutex pendant l'attente)
    \item \texttt{pthread\_cond\_signal()} : Réveille UN thread en attente
    \item \texttt{pthread\_cond\_broadcast()} : Réveille TOUS les threads en attente
    \item \texttt{pthread\_cond\_destroy()} : Détruit la variable de condition
\end{itemize}

\subsection{Exemple 6 : Producteur-Consommateur}

\begin{lstlisting}[caption={Problème classique producteur-consommateur}]
#include <stdio.h>
#include <stdlib.h>
#include <pthread.h>
#include <unistd.h>

#define TAILLE_BUFFER 5

(*@\color{commentcolor}// Buffer partagé@*)
int buffer[TAILLE_BUFFER];
int compteur = 0;  (*@\color{commentcolor}// Nombre d'éléments dans le buffer@*)

(*@\color{commentcolor}// Mutex et variables de condition@*)
pthread_mutex_t mutex = PTHREAD_MUTEX_INITIALIZER;
pthread_cond_t buffer_non_plein = PTHREAD_COND_INITIALIZER;
pthread_cond_t buffer_non_vide = PTHREAD_COND_INITIALIZER;

(*@\color{commentcolor}// Fonction producteur@*)
void* producteur(void* arg) {
    int id = *(int*)arg;
    
    for (int i = 0; i < 10; i++) {
        (*@\color{commentcolor}// Production d'un élément@*)
        int item = rand() % 100;
        
        (*@\color{commentcolor}// Verrouillage du mutex@*)
        pthread_mutex_lock(&mutex);
        
        (*@\color{commentcolor}// Attendre que le buffer ne soit pas plein@*)
        while (compteur == TAILLE_BUFFER) {
            printf("Producteur %d : Buffer plein, attente...\n", id);
            (*@\color{commentcolor}// pthread\_cond\_wait libère le mutex et attend@*)
            (*@\color{commentcolor}// Quand réveillé, il re-verrouille automatiquement le mutex@*)
            pthread_cond_wait(&buffer_non_plein, &mutex);
        }
        
        (*@\color{commentcolor}// Ajout de l'élément dans le buffer@*)
        buffer[compteur] = item;
        compteur++;
        
        printf("Producteur %d : Produit %d (buffer: %d/%d)\n",
               id, item, compteur, TAILLE_BUFFER);
        
        (*@\color{commentcolor}// Signaler qu'il y a un nouvel élément@*)
        pthread_cond_signal(&buffer_non_vide);
        
        (*@\color{commentcolor}// Déverrouillage du mutex@*)
        pthread_mutex_unlock(&mutex);
        
        (*@\color{commentcolor}// Pause aléatoire@*)
        usleep(rand() % 500000);
    }
    
    return NULL;
}

(*@\color{commentcolor}// Fonction consommateur@*)
void* consommateur(void* arg) {
    int id = *(int*)arg;
    
    for (int i = 0; i < 10; i++) {
        pthread_mutex_lock(&mutex);
        
        (*@\color{commentcolor}// Attendre que le buffer ne soit pas vide@*)
        while (compteur == 0) {
            printf("Consommateur %d : Buffer vide, attente...\n", id);
            pthread_cond_wait(&buffer_non_vide, &mutex);
        }
        
        (*@\color{commentcolor}// Retrait d'un élément du buffer@*)
        compteur--;
        int item = buffer[compteur];
        
        printf("Consommateur %d : Consommé %d (buffer: %d/%d)\n",
               id, item, compteur, TAILLE_BUFFER);
        
        (*@\color{commentcolor}// Signaler qu'il y a de la place@*)
        pthread_cond_signal(&buffer_non_plein);
        
        pthread_mutex_unlock(&mutex);
        
        usleep(rand() % 500000);
    }
    
    return NULL;
}

int main() {
    pthread_t prod[2], cons[2];
    int ids[2] = {1, 2};
    
    (*@\color{commentcolor}// Création de 2 producteurs et 2 consommateurs@*)
    for (int i = 0; i < 2; i++) {
        pthread_create(&prod[i], NULL, producteur, &ids[i]);
        pthread_create(&cons[i], NULL, consommateur, &ids[i]);
    }
    
    (*@\color{commentcolor}// Attente de tous les threads@*)
    for (int i = 0; i < 2; i++) {
        pthread_join(prod[i], NULL);
        pthread_join(cons[i], NULL);
    }
    
    (*@\color{commentcolor}// Nettoyage@*)
    pthread_mutex_destroy(&mutex);
    pthread_cond_destroy(&buffer_non_plein);
    pthread_cond_destroy(&buffer_non_vide);
    
    return 0;
}
\end{lstlisting}

% ============================================
\section{Attributs des Threads}
% ============================================

\subsection{Initialisation et Destruction}

\begin{lstlisting}[caption={Utilisation des attributs}]
#include <pthread.h>

pthread_attr_t attr;

(*@\color{commentcolor}// Initialisation des attributs@*)
pthread_attr_init(&attr);

(*@\color{commentcolor}// Configuration des attributs@*)
pthread_attr_setdetachstate(&attr, PTHREAD_CREATE_DETACHED);

(*@\color{commentcolor}// Création du thread avec attributs@*)
pthread_t thread;
pthread_create(&thread, &attr, fonction, NULL);

(*@\color{commentcolor}// Destruction des attributs@*)
pthread_attr_destroy(&attr);
\end{lstlisting}

\subsection{Attributs Principaux}

\begin{itemize}
    \item \textbf{Detach State} : \texttt{PTHREAD\_CREATE\_JOINABLE} ou \texttt{PTHREAD\_CREATE\_DETACHED}
    \item \textbf{Stack Size} : Taille de la pile du thread
    \item \textbf{Scheduling Policy} : Politique d'ordonnancement
    \item \textbf{Priority} : Priorité du thread
\end{itemize}

% ============================================
\section{Sémaphores POSIX}
% ============================================

\begin{definitionbox}{Qu'est-ce qu'un Sémaphore ?}
Un \textbf{sémaphore} est un compteur qui contrôle l'accès à une ressource partagée. Il peut autoriser plusieurs threads simultanés (contrairement au mutex).
\end{definitionbox}

\subsection{Fonctions Principales}

\begin{lstlisting}[caption={Utilisation des sémaphores}]
#include <semaphore.h>

sem_t semaphore;

(*@\color{commentcolor}// Initialisation : sem\_init(sem, pshared, valeur\_initiale)@*)
(*@\color{commentcolor}// pshared = 0 : partagé entre threads du même processus@*)
(*@\color{commentcolor}// valeur\_initiale : nombre de ressources disponibles@*)
sem_init(&semaphore, 0, 3);  (*@\color{commentcolor}// 3 ressources disponibles@*)

(*@\color{commentcolor}// Attendre (décrémente le compteur, bloque si 0)@*)
sem_wait(&semaphore);

(*@\color{commentcolor}// Section critique@*)
(*@\color{commentcolor}// ...@*)

(*@\color{commentcolor}// Signaler (incrémente le compteur)@*)
sem_post(&semaphore);

(*@\color{commentcolor}// Destruction@*)
sem_destroy(&semaphore);
\end{lstlisting}

\subsection{Exemple 7 : Pool de Connexions}

\begin{lstlisting}[caption={Limitation du nombre de connexions simultanées}]
#include <stdio.h>
#include <pthread.h>
#include <semaphore.h>
#include <unistd.h>

#define MAX_CONNEXIONS 3
#define NB_CLIENTS 10

(*@\color{commentcolor}// Sémaphore limitant à 3 connexions simultanées@*)
sem_t pool_connexions;

void* client(void* arg) {
    int id = *(int*)arg;
    
    printf("Client %d : Demande de connexion...\n", id);
    
    (*@\color{commentcolor}// Attendre qu'une connexion soit disponible@*)
    sem_wait(&pool_connexions);
    
    printf("Client %d : CONNECTÉ\n", id);
    
    (*@\color{commentcolor}// Simulation d'utilisation de la connexion@*)
    sleep(2);
    
    printf("Client %d : Déconnexion\n", id);
    
    (*@\color{commentcolor}// Libérer la connexion@*)
    sem_post(&pool_connexions);
    
    free(arg);
    return NULL;
}

int main() {
    pthread_t threads[NB_CLIENTS];
    
    (*@\color{commentcolor}// Initialisation : maximum 3 connexions simultanées@*)
    sem_init(&pool_connexions, 0, MAX_CONNEXIONS);
    
    (*@\color{commentcolor}// Création de 10 clients@*)
    for (int i = 0; i < NB_CLIENTS; i++) {
        int* id = malloc(sizeof(int));
        *id = i + 1;
        pthread_create(&threads[i], NULL, client, id);
        usleep(100000);  (*@\color{commentcolor}// Petit délai entre les créations@*)
    }
    
    (*@\color{commentcolor}// Attente de tous les clients@*)
    for (int i = 0; i < NB_CLIENTS; i++) {
        pthread_join(threads[i], NULL);
    }
    
    sem_destroy(&pool_connexions);
    
    return 0;
}
\end{lstlisting}

% ============================================
\section{Barrières de Synchronisation}
% ============================================

\begin{definitionbox}{Qu'est-ce qu'une Barrière ?}
Une \textbf{barrière} est un point de synchronisation où tous les threads doivent arriver avant que l'un d'eux puisse continuer.
\end{definitionbox}

\subsection{Exemple 8 : Calcul Parallèle par Phases}

\begin{lstlisting}[caption={Synchronisation avec barrière}]
#include <stdio.h>
#include <pthread.h>
#include <unistd.h>

#define NB_THREADS 4

(*@\color{commentcolor}// Barrière pour synchroniser les threads@*)
pthread_barrier_t barriere;

void* travailleur(void* arg) {
    int id = *(int*)arg;
    
    (*@\color{commentcolor}// Phase 1@*)
    printf("Thread %d : Phase 1 en cours...\n", id);
    sleep(id);  (*@\color{commentcolor}// Temps variable selon le thread@*)
    printf("Thread %d : Phase 1 terminée\n", id);
    
    (*@\color{commentcolor}// Attendre que tous les threads finissent la phase 1@*)
    pthread_barrier_wait(&barriere);
    
    (*@\color{commentcolor}// Phase 2 (commence seulement quand tous ont fini phase 1)@*)
    printf("Thread %d : Phase 2 en cours...\n", id);
    sleep(1);
    printf("Thread %d : Phase 2 terminée\n", id);
    
    (*@\color{commentcolor}// Attendre que tous finissent la phase 2@*)
    pthread_barrier_wait(&barriere);
    
    printf("Thread %d : Toutes les phases terminées\n", id);
    
    free(arg);
    return NULL;
}

int main() {
    pthread_t threads[NB_THREADS];
    
    (*@\color{commentcolor}// Initialisation de la barrière pour 4 threads@*)
    pthread_barrier_init(&barriere, NULL, NB_THREADS);
    
    for (int i = 0; i < NB_THREADS; i++) {
        int* id = malloc(sizeof(int));
        *id = i + 1;
        pthread_create(&threads[i], NULL, travailleur, id);
    }
    
    for (int i = 0; i < NB_THREADS; i++) {
        pthread_join(threads[i], NULL);
    }
    
    pthread_barrier_destroy(&barriere);
    
    return 0;
}
\end{lstlisting}

% ============================================
\section{Thread-Local Storage (TLS)}
% ============================================

\begin{definitionbox}{Variables Locales aux Threads}
Le \textbf{TLS} permet à chaque thread d'avoir sa propre copie d'une variable globale.
\end{definitionbox}

\subsection{Utilisation avec \_\_thread}

\begin{lstlisting}[caption={Variables thread-local}]
#include <stdio.h>
#include <pthread.h>

(*@\color{commentcolor}// Variable globale normale (partagée)@*)
int compteur_global = 0;

(*@\color{commentcolor}// Variable thread-local (chaque thread a sa copie)@*)
__thread int compteur_local = 0;

void* fonction(void* arg) {
    int id = *(int*)arg;
    
    for (int i = 0; i < 5; i++) {
        compteur_global++;  (*@\color{commentcolor}// Partagé entre tous@*)
        compteur_local++;   (*@\color{commentcolor}// Propre à ce thread@*)
        
        printf("Thread %d : Global=%d, Local=%d\n",
               id, compteur_global, compteur_local);
    }
    
    return NULL;
}

int main() {
    pthread_t t1, t2;
    int id1 = 1, id2 = 2;
    
    pthread_create(&t1, NULL, fonction, &id1);
    pthread_create(&t2, NULL, fonction, &id2);
    
    pthread_join(t1, NULL);
    pthread_join(t2, NULL);
    
    return 0;
}
\end{lstlisting}

% ============================================
\section{Gestion des Erreurs}
% ============================================

\subsection{Codes d'Erreur Courants}

\begin{center}
\begin{tabular}{|l|p{8cm}|}
\hline
\textbf{Code} & \textbf{Signification} \\
\hline
EAGAIN & Ressources système insuffisantes \\
\hline
EINVAL & Argument invalide \\
\hline
EDEADLK & Deadlock détecté \\
\hline
EBUSY & Ressource occupée \\
\hline
EPERM & Opération non permise \\
\hline
\end{tabular}
\end{center}

\subsection{Exemple de Gestion d'Erreurs}

\begin{lstlisting}[caption={Vérification des erreurs}]
#include <stdio.h>
#include <pthread.h>
#include <errno.h>
#include <string.h>

void* fonction(void* arg) {
    return NULL;
}

int main() {
    pthread_t thread;
    int ret;
    
    (*@\color{commentcolor}// Création du thread avec vérification@*)
    ret = pthread_create(&thread, NULL, fonction, NULL);
    if (ret != 0) {
        fprintf(stderr, "Erreur pthread_create: %s\n", strerror(ret));
        return 1;
    }
    
    (*@\color{commentcolor}// Join avec vérification@*)
    ret = pthread_join(thread, NULL);
    if (ret != 0) {
        fprintf(stderr, "Erreur pthread_join: %s\n", strerror(ret));
        return 1;
    }
    
    return 0;
}
\end{lstlisting}

% ============================================
\section{Problèmes Classiques et Solutions}
% ============================================

\subsection{Deadlock (Interblocage)}

\begin{importantbox}{Conditions du Deadlock}
Un deadlock se produit quand :
\begin{enumerate}
    \item Exclusion mutuelle
    \item Hold and wait (détenir et attendre)
    \item Pas de préemption
    \item Attente circulaire
\end{enumerate}
\end{importantbox}

\subsection{Exemple de Deadlock}

\begin{lstlisting}[caption={Situation de deadlock}]
pthread_mutex_t mutex1 = PTHREAD_MUTEX_INITIALIZER;
pthread_mutex_t mutex2 = PTHREAD_MUTEX_INITIALIZER;

void* thread1_func(void* arg) {
    pthread_mutex_lock(&mutex1);    (*@\color{commentcolor}// Thread 1 verrouille mutex1@*)
    sleep(1);
    pthread_mutex_lock(&mutex2);    (*@\color{importantcolor}// Attend mutex2 (détenu par thread2)@*)
    
    (*@\color{commentcolor}// Section critique@*)
    
    pthread_mutex_unlock(&mutex2);
    pthread_mutex_unlock(&mutex1);
    return NULL;
}

void* thread2_func(void* arg) {
    pthread_mutex_lock(&mutex2);    (*@\color{commentcolor}// Thread 2 verrouille mutex2@*)
    sleep(1);
    pthread_mutex_lock(&mutex1);    (*@\color{importantcolor}// Attend mutex1 (détenu par thread1)@*)
    
    (*@\color{commentcolor}// Section critique@*)
    
    pthread_mutex_unlock(&mutex1);
    pthread_mutex_unlock(&mutex2);
    return NULL;
}
\end{lstlisting}

\subsection{Solution : Ordre de Verrouillage}

\begin{lstlisting}[caption={Éviter le deadlock avec un ordre fixe}]
void* thread1_func(void* arg) {
    (*@\color{commentcolor}// Toujours verrouiller dans le même ordre@*)
    pthread_mutex_lock(&mutex1);
    pthread_mutex_lock(&mutex2);
    
    (*@\color{commentcolor}// Section critique@*)
    
    pthread_mutex_unlock(&mutex2);
    pthread_mutex_unlock(&mutex1);
    return NULL;
}

void* thread2_func(void* arg) {
    (*@\color{commentcolor}// Même ordre de verrouillage@*)
    pthread_mutex_lock(&mutex1);
    pthread_mutex_lock(&mutex2);
    
    (*@\color{commentcolor}// Section critique@*)
    
    pthread_mutex_unlock(&mutex2);
    pthread_mutex_unlock(&mutex1);
    return NULL;
}
\end{lstlisting}

% ============================================
\section{Bonnes Pratiques}
% ============================================

\begin{examplebox}{Règles d'Or}
\begin{enumerate}
    \item \textbf{Toujours} compiler avec \texttt{-pthread}
    \item \textbf{Toujours} vérifier les valeurs de retour
    \item \textbf{Minimiser} la taille des sections critiques
    \item \textbf{Éviter} les variables globales quand possible
    \item \textbf{Utiliser} des mutex pour protéger les données partagées
    \item \textbf{Libérer} toujours les mutex verrouillés
    \item \textbf{Détruire} les ressources (mutex, threads, etc.)
    \item \textbf{Éviter} les deadlocks avec un ordre de verrouillage cohérent
    \item \textbf{Utiliser} \texttt{pthread\_join()} ou \texttt{pthread\_detach()}
    \item \textbf{Tester} avec différents nombres de threads
\end{enumerate}
\end{examplebox}

% ============================================
\section{Résumé des Fonctions Principales}
% ============================================

\subsection{Tableau Récapitulatif}

\begin{center}
\small
\begin{tabular}{|l|p{9cm}|}
\hline
\textbf{Fonction} & \textbf{Description} \\
\hline
\multicolumn{2}{|c|}{\textbf{Gestion des Threads}} \\
\hline
pthread\_create() & Crée un nouveau thread \\
pthread\_join() & Attend la terminaison d'un thread \\
pthread\_exit() & Termine le thread appelant \\
pthread\_detach() & Détache un thread \\
pthread\_self() & Retourne l'ID du thread appelant \\
pthread\_equal() & Compare deux IDs de threads \\
\hline
\multicolumn{2}{|c|}{\textbf{Mutex}} \\
\hline
pthread\_mutex\_init() & Initialise un mutex \\
pthread\_mutex\_lock() & Verrouille un mutex \\
pthread\_mutex\_unlock() & Déverrouille un mutex \\
pthread\_mutex\_trylock() & Tente de verrouiller sans bloquer \\
pthread\_mutex\_destroy() & Détruit un mutex \\
\hline
\multicolumn{2}{|c|}{\textbf{Variables de Condition}} \\
\hline
pthread\_cond\_init() & Initialise une variable de condition \\
pthread\_cond\_wait() & Attend un signal \\
pthread\_cond\_signal() & Réveille un thread \\
pthread\_cond\_broadcast() & Réveille tous les threads \\
pthread\_cond\_destroy() & Détruit une variable de condition \\
\hline
\multicolumn{2}{|c|}{\textbf{Sémaphores}} \\
\hline
sem\_init() & Initialise un sémaphore \\
sem\_wait() & Décrémente (P operation) \\
sem\_post() & Incrémente (V operation) \\
sem\_destroy() & Détruit un sémaphore \\
\hline
\multicolumn{2}{|c|}{\textbf{Barrières}} \\
\hline
pthread\_barrier\_init() & Initialise une barrière \\
pthread\_barrier\_wait() & Attend à la barrière \\
pthread\_barrier\_destroy() & Détruit une barrière \\
\hline
\end{tabular}
\end{center}

% ============================================
\section{Exercices Types pour l'Examen}
% ============================================

\subsection{Exercice 1 : Analyse de Code}

\begin{examplebox}{Question}
Identifiez le problème dans ce code et proposez une correction :
\end{examplebox}

\begin{lstlisting}
int compteur = 0;

void* incrementer(void* arg) {
    for (int i = 0; i < 100000; i++) {
        compteur++;
    }
    return NULL;
}

int main() {
    pthread_t t1, t2;
    pthread_create(&t1, NULL, incrementer, NULL);
    pthread_create(&t2, NULL, incrementer, NULL);
    pthread_join(t1, NULL);
    pthread_join(t2, NULL);
    printf("Compteur: %d\n", compteur);
    return 0;
}
\end{lstlisting}

\textbf{Réponse :} Race condition sur \texttt{compteur}. Solution : utiliser un mutex.

\subsection{Exercice 2 : Implémentation}

\begin{examplebox}{Question}
Écrivez un programme avec 3 threads qui calculent la somme de parties différentes d'un tableau, puis combinent les résultats.
\end{examplebox}

\subsection{Exercice 3 : Synchronisation}

\begin{examplebox}{Question}
Implémentez le problème des lecteurs-rédacteurs : plusieurs lecteurs peuvent lire simultanément, mais un seul rédacteur à la fois, et pas de lecteur pendant l'écriture.
\end{examplebox}

% ============================================
\section{Conclusion}
% ============================================

\begin{importantbox}{Points Clés à Retenir}
\begin{itemize}
    \item Les threads partagent la mémoire $\Rightarrow$ besoin de synchronisation
    \item Mutex = exclusion mutuelle (1 seul thread à la fois)
    \item Variables de condition = attente d'événements
    \item Sémaphores = compteur de ressources
    \item Toujours protéger les sections critiques
    \item Attention aux deadlocks !
    \item Compiler avec \texttt{-pthread}
\end{itemize}
\end{importantbox}

\vfill

\begin{center}
\Large\textbf{Bonne chance pour votre examen !}
\end{center}

\end{document}
